\section*{Qué queda al restar el aporte de cada cuadrado}
Fijemos la máscara basal binaria \(M\) y las incidencias diagonales físicas.
Para \(Y\) contenido en sus sitios, actívense sólo las incidencias entre sitios de \(Y\).
Siempre se mantienen las restricciones de torres y el exterior de la misma máscara.
Sea \(e(Y)\) el logaritmo de la probabilidad de evitación bajo torres menos
el cambio del predictor físico respecto al de torres. Sea
\[
 e[U]=\sum_{Y\subseteq U}(-1)^{|U|-|Y|}e(Y).
\]
En la llegada corta las positividades y el reemplazo del aproximante tienen
los dominios de v3/06--07. En particular \(e(\varnothing)=0\) y \(e(M)=Y_D\).
Particionemos los pares físicos por su cuadrado geométrico \(C\). Éste se determina
por sus dos filas y sus dos columnas, incluso si sólo un par está disponible.
Sea \(e_C(Y)\) el mismo error activando sólo los pares de \(C\) contenidos en \(Y\).
En el conjunto pleno \(e_C(M)=\gamma_C\), por la definición de la corrección exacta.
Así \(g^\flat=e(M)-\sum_C e_C(M)\). Esta proyección opera sobre las activaciones,
sin formar el menor de torres al apagar un sitio.

Si \(U\) contiene un sitio ajeno a todo par activable en \(U\), la alternancia
se anula por emparejamiento de las caras que sólo difieren en ese sitio.
Si todos los pares de \(U\) pertenecen al mismo cuadrado, \(e(Y)=e_C(Y)\) en
todas sus caras y los demás \(e_{C'}(Y)\) son cero. Si hay al menos dos cuadrados,
ningún \(e_C[U]\) puede contribuir salvo si todos los sitios son vértices de \(C\).
En ese último caso las únicas parejas diagonales posibles son las dos de \(C\).
Por tanto el coeficiente restante es el crudo \(e[U]\) sobre soportes con al menos
dos cuadrados distintos y sin sitios aislados. En particular, el grado de la
sustracción es a lo sumo cuatro.

\paragraph{Resto físico en los tres sitios compatibles.}
Con \(z,a,b\) en filas y columnas distintas, sean \(A,B\) las probabilidades
exactas de los pares \(za,zb\) e \(I\) la de los tres sitios, en la medida de torres.
En B1 la respuesta del conteo es
\(\log(1+(I-AB)/((1-A)(1-B)))= I-AB+O(j^{-5})\).
En C1 hay tres pares en una línea. La respuesta del conteo tras restar las
tres respuestas de par es \(2I-\sum_{P<P'}\pi^{\rm ex}_P\pi^{\rm ex}_{P'}+O(j^{-5})\).
En ambos casos el grafo con sus incidencias y fibras de torres es un árbol con
tres sitios. V3/05 (apartados de cuñas y de grado bajo) aplica P05 antes de
las sustracciones. Con una componente da \(p=3+1+1=5\).
En B1 la alternancia cruda resta los dos pares. En C1 resta los tres.
Así las dos respuestas físicas del predictor tienen los mismos términos cúbicos
y cuárticos, con resto uniforme \(O(j^{-5})\) en la clase de márgenes fijados.
Los polinomios explícitos son los de la parte 08 de la entrega anterior.
Este uso de P05 conserva el cofactor y el sitio que pertenece a dos líneas en B1.
Una coincidencia adicional de fila o columna entre los extremos cambia el grafo
y se rige por el exponente de ciclo de v3/05.
